% ==== partgf2 论文补充部分:第一/二/三讲的补充与参考文献 ====
% 源: MinerU_latex_A-非.tex 4615-6327 行(GEOMETRIC AND FOURIER METHODS FOR NONLINEAR
%     WAVE EQUATIONS 的第 4、5、6 节与 REFERENCES;接 partgf1, 公式编号自 (94) 起)
% 术语: 共形等距=conformal isometry; 等距=isometry; 规范条件=gauge condition;
%       零无穷大=null infinity; 未来零无穷大=future null infinity; 零测地生成线=null geodesic generator;
%       仿射参数=affine parameter; 影响域=domain of influence; 因果未来=causal future;
%       电磁场=electromagnetic field; 无散度=divergence-free; 能量-动量张量=energy-momentum tensor;
%       共形守恒律=conformal conservation law; 广义能量估计=generalized energy estimates;
%       零二次型=null quadratic form; 齐次 Sobolev 空间=homogeneous Sobolev space;
%       频率依赖区间=frequency dependent intervals; 重标度=rescaling; 色散不等式=dispersive inequality;
%       弱/强宇宙监督猜想=weak/strong cosmic censorship conjecture; 时间规范=Temporal gauge
% REFS:
%   - "Theorem 1.2"(定义 6.2 中 "the set Ω0 of Theorem 1.2"):论文第 1 讲无定理 1.2
%     (第 2 个编号对象为注 1.2), 疑为原书笔误(应为定理 6.1 或主书第 1 章定理 1.2), 暂按
%     "定理 1.2" 直译, 待统一解析. 
%   - "so (119) is invariant"(源 5542 行)与 "(119) takes the form"(5546 行):原书把显示为
%     (YM) 的方程 D^νF_μν=0 称为 (119), 而 (119) 又是 Coulomb 规范方程组 (119a-c) 的编号, 
%     原书自身不一致;两处按原文直译 "(119)", 不做链接. 
%   - "the wave equation (1)"(4624 行)→ \eqref{eq:cG-1};"(31)"/"(32)"(5114、5186 行)→
%     \eqref{eq:cG-31}/\eqref{eq:cG-32};"(67)"(5705、5728、5884 行)→ \eqref{eq:cG-67};
%     "(71)"(5721 行)→ \eqref{eq:cG-71};(均为 partgf1 内标签)
%   - 定理 3.10/3.12(5690-5693 行)→ \thmref{thm:cG-3.10}/\thmref{thm:cG-3.12}(partgf1);
%   - FOOTNOTE: 论文部分在 pdf_text.txt 中为乱码, 脚注 29-44 原文从作者论文
%     (IPAM 版 oide2001_4284.pdf) 找回并翻译;脚注 32 中两处引文原书即为 "[?]"(作者未填), 照译保留. 
%   - OCR-BROKEN/存疑(均按原书 PDF 校过):
%     * "Observe that C has the important property..."(4657 行):"C" 在原书中未定义, 照直译并待核;
%     * 命题 4.10 中 "Ω = 4(n+1)t"(4914 行):两版原书均如此, 照译(疑应为 4t);
%     * 定义 4.12 中 \widetilde{\nabla}/\vec{\nabla}/\nabla 均为 OCR 变体, 统一为 \nabla;
%     * (96) 右端按数学含义与 (94)/(95) 对称性还原为 \underline{L}(φ)²;
%     * Q_0(φ, ψ) 展开(5170 行)第二项 OCR 误为 E_+φ·E_-ψ, 已按原书校为 E_-φ·E_+ψ;
%     * Q_ij(φ, ψ) 展开(5183 行)末项按原书 PDF 字符流还原为 ∂_rφ((x_i/r)A_jφ − (x_j/r)A_iφ), 
%       原书该式各因子均含 φ, 疑为笔误(数学上应为 ψ), 照录待核;
%     * 引理 4.16 证明中 "ξ ∈ R^{n-1}" 原书如此(疑应为单位球面 S^{n-1}), 照录;
%     * 命题 4.5 后补回 OCR 丢失的 "if" 部分证明段(据原书 PDF);
%     * (116) 前 "-Δ□A_j"(5433 行)保留负号(ipam 版缺负号, 数学上应有);
%     * "whichimplies"(5466 行)按 "which implies" 处理;源 5755 行 "R^#/{V}" 还原为 R^3\{0};
%     * 源 5648 行 "D_i^+" 还原为 D_0^+;源 5802 行 "S_λ^aφ" 还原为 S_{λ^a}φ;
%     * 源 5408 行 "\lceil(∂_⊔A_I)" 按原书还原为 d(∂_tA_0);源 4790 行 K_μ 公式补括号. 
% 编号说明:
%   - 公式 (94)-(147) 按原书连续编号(接 partgf1 的 (93));字母编号公式 (MKGa)/(MKGb)/(YM)
%     沿用 partgf1 的 \tag 处理(指南"禁止 \tag"规则的特例, 否则后续编号整体偏移);
%     (110a-c)、(111a-b)、(113a-d)、(114a-b)、(119a-c) 用 subequations, 每组消耗一个编号;
%   - 第 4 节:4.1、4.8、4.14 为共享序列子标题(\mtsubsec), 与定义/命题/定理/引理/注共享计数器, 
%     自然编号 4.2-4.17;
%   - 第 5 节:仅 5.0.1/5.0.2/5.0.3 三个子标题(原书无编号定理对象), 用 \mtsubsecn 保持
%     原显示编号且不推进计数器;
%   - 第 6 节:6.3、6.4、6.6、6.8、6.11、6.13、6.16、6.18 为 \mtsubsec 子标题;
%     Remark (A2)、(A3) 与 6.14 后的 Remark 未编号, 用 remark*;两处 Conjecture 用加粗标题;
%     定理 (A)=定理 3.10、(B)=定理 3.12(partgf1);定理 6.7/6.9/6.12/6.15/6.17/6.19/6.20
%     按任务要求以 (A1)-(A5)、(A6)、(B) 为环境标题;
%   - 参考文献 1-76 条:[32] 末尾 "46 (1993),  1221--1268" 与 [30] 中的重复为原书残留, 照录. 
% ====

\subsection{第一讲的补充}

像 Yang-Mills 方程、波映射与 Einstein 方程这样重要的非线性波动方程, 都具有与其几何特性密切相关的根本对称性. 深入研究微分几何, 尤其是 Lorentz 几何, 在研究这些方程时至少与调和分析同等重要. 下面我们将考察 Minkowski 时空的共形结构与波动方程 \eqref{eq:cG-1} 的估计之间的关系. 特别是, 我们将使用下面第三讲中所讨论的广义能量估计. 我假定读者对微分几何的基本概念有一定了解, 如 Lie 导数、协变导数以及曲率张量. 

\mtsubsec{Minkowski 空间的共形结构}

\begin{definition}\label{def:cG-4.2}
考虑一个时空 $(\mathcal{M},  g)$, 即一个带 Lorentz 度量 $g$ 的 Lorentz 流形, 其符号为 $(-1,  1,  \ldots,  1)$. 

微分同胚 $\Phi: \mathcal{U} \subset \mathcal{M} \to \mathcal{M}$ 称为共形等距, 若在每一点 $p$ 处 $\Phi_{*}g = \Lambda^{2}g$, 即

\[
\left(\Phi^ {*} g\right) (X,  Y) | _ {p} = g \left(\Phi_ {*} X,  \Phi_ {*} Y\right) | _ {\Phi (p)} = \Lambda^ {2} g (X,  Y) | _ {p}
\]

其中 $\Lambda \neq 0$. 若 $\Lambda = 1$, 则称 $\Phi$ 为 $\mathcal{M}$ 的等距. 
\end{definition}

\begin{definition}\label{def:cG-4.3}
生成单参数等距群(分别地, 共形等距群)的向量场 $K$ 称为 Killing(分别地, 共形 Killing)向量场. 
\end{definition}

设 $K$ 是这样一个向量场, $\Phi_{t}$ 是对应的单参数群. 由于 $(\Phi_{t})_{*}$ 是共形等距, 我们推知 $L_{K}g$ 必定与度量 $g$ 成比例. 此外, 若 $K$ 是 Killing 向量场, 则 $L_{K}g = 0$. 

注意, $C$ 具有关于共形等距为共形不变的重要性质. 

\begin{definition}\label{def:cG-4.4}
给定任意向量场 $X$, 记 $^{(X)}\pi$ 为 $X$ 的形变张量, 由如下公式定义

\[
{ } ^ { ( X ) } \pi _ { \alpha \beta } = ( \mathcal { L } _ { X } g ) _ { \alpha \beta } = D _ { \alpha } X _ { \beta } + D _ { \beta } X _ { \alpha } .
\]
\end{definition}

张量 $^{(X)}\pi$ 在精确的意义上度量了由 $X$ 生成的微分同胚与等距或共形等距相差多少. 

\begin{proposition}\label{prop:cG-4.5}
向量场 $X$ 是 Killing 的当且仅当 $^{(X)}\pi = 0$. 它是共形 Killing 的当且仅当 $^{(X)}\pi$ 与 $g$ 成比例. 
\end{proposition}

为证明"若"部分, 我们可以选取局部坐标 $x^{0},  x^{1},  \ldots,  x^{n}$ 使得 $X = \frac{\partial}{\partial x^{0}}$. 于是立即推出, 相对于这些坐标, 度量 $g$ 与 $x^{0}$ 无关. 

\begin{proposition}\label{prop:cG-4.6}
在任何时空 $\mathcal{M}$ 上, 线性无关的 Killing 向量场不可能多于 $\frac{1}{2} (n + 1)(n + 2)$ 个. 
\end{proposition}

具有最多 Killing 与共形 Killing 向量场的时空是 Minkowski 时空. 让我们回顾与之相关的等距与共形等距. 

设 $x^{\mu}$ 是一个正向的惯性坐标系, 我们有:

\begin{enumerate}
\item 平移:对任意给定向量 $a = (a^{0},  a^{1},  \ldots,  a^{n}) \in \mathbf{M}$, 
\[
x^{\mu} \rightarrow x^{\mu} + a^{\mu}
\]
\item Lorentz 旋转:给定任意 $\Lambda = \Lambda_{\sigma}^{\rho} \in \mathbf{O}(1,  n)$, 
\[
x^{\mu} \rightarrow \Lambda_{\nu}^{\mu} x^{\nu}
\]
\item 伸缩:给定任意实数 $\lambda \neq 0$, 
\[
x^{\mu} \to \lambda x^{\mu}
\]
\item 反演:考虑变换 $x^{\mu} \rightarrow I(x^{\mu})$, 其中
\[
I (x^{\mu}) = \frac {x^{\mu}}{(x ,  x)}
\]
对所有满足 $(x,  x) \neq 0$ 的点 $x \in M$ 定义. 
\end{enumerate}

前两个变换是 $M$ 的等距, 由它们生成的群称为 Poincar\'e 群. 后两类变换是共形等距. 由以上所有变换生成的群称为共形群. 让我们列出生成上述变换的 Killing 与共形 Killing 向量场. 

\begin{enumerate}
\renewcommand{\labelenumi}{\roman{enumi}.}
\item $x^{\mu}$ 方向上的平移生成元, $\mu = 0,  1,  \ldots,  n$:
\[
T_{\mu} = \frac{\partial}{\partial x^{\mu}}
\]
\item $(\mu,  \nu)$ 平面上的 Lorentz 旋转生成元:
\[
L_{\mu\nu} = x_{\mu}\partial_{\nu} - x_{\nu}\partial_{\mu}
\]
\item 伸缩变换的生成元:
\[
S = x^{\mu}\partial_{\mu}
\]
\item 倒置平移的生成元:
\[
K_{\mu} = 2x_{\mu}(x^{\rho}\partial_{\rho}) - (x^{\rho}x_{\rho})\partial_{\mu}
\]
\end{enumerate}

\begin{theorem}\label{thm:cG-4.7}
对 $n > 1$, $\mathbf{M}^{n + 1}$ 中任何共形 Killing 向量场都是向量场 $T,  L,  S$ 与 $K$ 的实常数系数的线性组合. 若 $n = 1$, 所有共形向量场由如下公式给出

\[
f (t + x) (\partial_ {t} + \partial_ {x}) + g (t - x) (\partial_ {t} - \partial_ {x})
\]

其中 $f,  g$ 是单变量的任意光滑函数. 
\end{theorem}
\mtsubsec{广义能量估计}

考虑平坦 Minkowski 空间 $\mathbb{R}^{n+1}$ 中的波动方程 $\Box\phi=0$ 及其相伴的能量-动量张量, 

\[
Q_{\alpha\beta} = Q_{\alpha\beta}[\phi] = \partial_{\alpha}\phi\, \partial_{\beta}\phi - \frac{1}{2}m_{\alpha\beta}\big(m^{\mu\nu}\partial_{\mu}\phi\, \partial_{\nu}\phi\big).
\]

\begin{proposition}\label{prop:cG-4.9}
能量-动量张量是对称且无散的, 即

\[
\partial^{\beta}Q_{\alpha\beta} = 0.
\]

此外, 对任意未来指向的类时向量场 $X,  Y$, 有 $Q(X,  Y) \geq 0$. 
\end{proposition}

\begin{proof}
前两个论断是显然的. 为检验第三个, 考虑任意一点处的零对\footnote{两个未来指向的线性无关的零向量 $L,  \underline{L}$, 满足 $\langle L,  L\rangle = \langle \underline{L},  \underline{L}\rangle = 0$ 以及归一化条件 $\langle L,  \underline{L}\rangle = 2$. }$L,  \underline{L}$, 它们张成一个包含 $X$ 与 $Y$ 的平面. 于是 $X,  Y$ 是 $L,  \underline{L}$ 的正系数线性组合. 正性于是是如下简单而重要计算的直接推论, 

\begin{equation}
Q ( L ,  L ) = L ( \phi ) ^ { 2 }\label{eq:cG-94}
\end{equation}

\begin{equation}
Q (L,  \underline {{L}}) = | \nabla \phi | ^ {2}\label{eq:cG-95}
\end{equation}

\begin{equation}
Q (\underline {{L}},  \underline {{L}}) = \underline {L} (\phi) ^ {2}\label{eq:cG-96}
\end{equation}

其中 $|\nabla \phi| = \sum_{A}|e_{A}(\phi)|^{2}$, $(e_A)_{A=1, \dots , n - 1}$ 是张成 $L,  \underline{L}$ 正交补的标准正交标架. 
\end{proof}

设 $X$ 是任意向量场, 考虑与其相伴的时空动量向量场 $P_{\alpha} = Q_{\alpha\beta}X^{\beta}$. 利用 $Q$ 的对称性与无散性以及形变张量 $\pi = {}^{(X)}\pi$ 的定义, 我们推知

\[
\partial_{\alpha}P^{\alpha} = \frac{1}{2}Q^{\alpha\beta}\pi_{\alpha\beta}.
\]

因此, 若 $X$ 是 Killing 向量场, 则 $\partial_{\alpha}P^{\alpha}=0$, 积分即得熟悉的守恒律. 特别地, 若我们在时空片 $[t_{0}, t]\times R^{n}$ 上积分, 则导出

\begin{equation}
\int_{\Sigma_{t}} Q(X,  \partial_{t}) dx = \int_{\Sigma_{t_{0}}} Q(X,  \partial_{t}) dx\label{eq:cG-97}
\end{equation}

其中 $\partial_{t}$ 在这里扮演类空超曲面 $\Sigma_{t}$ 的未来指向单位法向量的角色. 取 $X$ 亦为 Killing 向量场 $\partial_{t}$, 我们导出熟悉的能量守恒律

\[
\int_{\Sigma_{t}} Q(\partial_{t},  \partial_{t}) dx = \int_{\Sigma_{t_{0}}} Q(\partial_{t},  \partial_{t}) dx
\]

其中 $Q(\partial_t,  \partial_t) = \frac{1}{2} \left( |\partial_t \phi|^2 + \sum_{i=1}^{n} |\partial_i \phi|^2 \right)$ 是通常的能量密度. 动量与角动量守恒律也包含在 \eqref{eq:cG-97} 中, 只要取 $X = \partial_i$ 或 $X = L_{ij}$. 

若我们不在片 $[0, T]\times \mathbb{R}^n$ 上积分恒等式 $\partial^{\alpha}P_{\alpha}$, 而是在未来区域\footnote{或依赖域}$\mathcal{D}$ 上积分, 其中

\[
\mathcal {D} = \{(t,  x) / | x - y | \leq T - t,  t _ {0} \leq t \leq t _ {1} \leq T \}, 
\]

则导出守恒律

\begin{equation}
\int_{\Sigma_{t_{1}}} Q(X,  \partial_{t}) dx + \int_{\underline{C}(t_{0},  t_{1})} Q(X,  \underline{L}) = \int_{\Sigma_{t_{0}}} Q(X,  \partial_{t}) dx\label{eq:cG-98}
\end{equation}

其中 $\underline{C}(t_{0}, t_{1})$ 是 $\mathcal{D}$ 的外边界, $\underline{L}=\partial_{t}-\sum_{i}\frac{(x^{i}-x_{0}^{i})}{|x-x_{0}|}\partial_{i}$ 是内零超曲面 $C$ 的零向量场法向量\footnote{在 Minkowski 时空度量的意义下}, 通量 $\int_{\underline{C}(t_{0}, t_{1})}Q(X, \underline{L})=\int_{t_{0}}^{t_{1}}dt\int_{S_{\underline{u}, t}}Q(X, \underline{L})$, 其中 $S_{\underline{u}, t}$ 是 $C$ 与 $\Sigma_{t}$ 的交球面. 在 $X=\partial_{t}$ 的特殊情形下, $\Sigma_{t}$ 上的两个能量积分与通量均为正, 这给出波动方程依赖域性质的一个简单证明. 

上述导致 \eqref{eq:cG-97} 的过程可以推广到共形 Killing 向量场. 

\begin{proposition}\label{prop:cG-4.10}
设 $Q_{\alpha\beta} = Q_{\alpha\beta}[\phi]$ 是 $\square\phi = 0$ 的解对应的能量-动量张量. 设 $X$ 是共形 Killing 向量场, 即 $\pi = {}^{(X)}\pi = \mathcal{L}_Xm = \Omega m$, 且 $tr\pi = m^{\alpha\beta}\pi_{\alpha\beta}$. 容易验证 $\square\Omega = 0$;事实上, 在 $X = K_0$ 的特殊情形下, $\Omega = 4(n + 1)t$. 设

\[
\bar {P} _ {\alpha} = Q _ {\alpha\beta} X^{\beta} + \frac {n - 1}{4 (n + 1)} tr\pi\, \phi\, \partial_{\alpha}\phi - \frac {n - 1}{8 (n + 1)}\partial_{\alpha}(tr\pi)\phi^{2}.
\]

我们有

\[
\partial^{\alpha}\bar {P} _ {\alpha} = 0.
\]
\end{proposition}

将该命题应用于共形 Killing 向量场 $X = K := K_{0}$, 并在时间片 $[t_{0},  t] \times R^{n}$ 上积分相应的无散方程, 我们推得如下结论\footnote{该命题的第 $i$ 部分与第 ii 部分归功于 C. Morawetz [?]. 第 iii 部分见 [?], 第 310--313 页. }:

\begin{proposition}\label{prop:cG-4.11}
设 $\overline{Q}(K_0, T_0) = Q(K_0, T_0) + (n - 1)t\phi\, \partial_t\phi -\frac{n - 1}{2}\phi^2$, 其中 $T_{0} = \partial_{t}$ 是 $\Sigma_t$ 的单位法向量, $\phi$ 是 $\square \phi = 0$ 的解. 

i.) 如下共形守恒律成立, 

\begin{equation}
\int_{\Sigma_{t}} \overline{Q}(K_{0},  T_{0}) = \int_{\Sigma_{t_{0}}} \overline{Q}(K_{0},  T_{0})\label{eq:cG-99}
\end{equation}

ii.) 此外我们有

\begin{equation}
\int_{\Sigma_{t}} \overline{Q}(K_{0},  T_{0}) = \frac{1}{4}\bigg(\int_{\Sigma_{t}} \underline{u}^{2}(L^{\prime}\phi)^{2} + \int_{\Sigma_{t}} 2(t^{2} + r^{2})|\nabla\phi|^{2} + \int_{\Sigma_{t}} u^{2}(\underline{L}^{\prime}\phi)^{2}\bigg)\label{eq:cG-100}
\end{equation}

其中 $L = \partial_t + \partial_r$, $\underline{L} = \partial_t - \partial_r$, $u = t - r$, $\underline{u} = t + r$, 且 $\underline{u}L'(\phi) = \underline{u}L(\phi) + (n - 1)\phi$, $u\underline{L}'(\phi) = u\underline{L}(\phi) + (n - 1)\phi$. 

iii.) 又若 $n \geq 3$, 则存在常数 $c > 0$ 使得

\begin{equation}
\int_{\Sigma_{t}} \overline{Q}(K_{0},  T_{0}) \geq c\left(\int_{\Sigma_{t}} \underline{u}^{2}(L\phi)^{2} + \int_{\Sigma_{t}} 2(t^{2} + r^{2})|\nabla\phi|^{2} + \int_{\Sigma_{t}} u^{2}(\underline{L}\phi)^{2} + \phi^{2}\right)\label{eq:cG-101}
\end{equation}
\end{proposition}

\begin{definition}\label{def:cG-4.12}
我们引入如下范数:

\[
\begin{array}{rcl}
E[\phi](t) & = & \int_{\Sigma_{t}}|\phi(t)|^{2} \\
\mathcal{Q}[\phi](t) & = & \int_{\Sigma_{t}}\overline{Q}[\phi](K_{0},  T_{0}) \\
\mathcal{E}[\phi](t) & = & \int_{\Sigma_{t}}\bigg(\underline{u}^{2}|L\phi|^{2} + 2(t^{2} + r^{2})|\nabla\phi|^{2} + u^{2}|\underline{L}\phi|^{2} + |\phi|^{2}\bigg)
\end{array}
\]

我们还引入高阶范数:

\begin{align*}
E_{s}[\phi](t) &= \sum_{j = 0}^{s}\sum_{\Gamma_{i_{1}}, ..,  \Gamma_{i_{j}}} E[\Gamma_{i_{1}}, ..,  \Gamma_{i_{j}}\phi](t) \\
\mathcal{Q}_{s}[\phi](t) &= \sum_{j = 0}^{s}\sum_{\Gamma_{i_{1}}, ..,  \Gamma_{i_{j}}} \mathcal{Q}[\Gamma_{i_{1}}, ..,  \Gamma_{i_{j}}\phi](t)
\end{align*}

其中 $\Gamma$ 是下列向量场中的任意一个

\[
S = t\partial_{t} + \sum_{i}x_{i}\partial_{i},  \qquad L_{i} = x_{i}\partial_{t} + t\partial_{i},  \qquad \Omega_{ij} = x_{i}\partial_{j} - x_{j}\partial_{i}.
\]
\end{definition}

\begin{remark}\label{rem:cG-4.13}
注意

\[
E_{s + 1}(t) \lesssim \mathcal{Q}_{s}(t) = \mathcal{Q}_{s}(0).
\]
\end{remark}
\mtsubsec{衰减估计}

众所周知, $\mathbb{R}^{n+1}$ 中自由波动方程 $\Box\phi = 0$ 的解在一致范数下当 $t \to \infty$ 时按 $t^{-\frac{n-1}{2}}$ 衰减. 此外, 在向外特征方向之外衰减更快. 这类信息可以从物理空间或 Fourier 空间中解的显式表示导出. 本小节的目的是说明, 如何不需要显式表示、仅用上一小节讨论的广义能量估计即可获得这类信息. 

考虑标准零对 $E_{\pm} = \partial_t \pm \partial_r$, 以及角向向量场 $A_i = \partial_i - \frac{x_i}{r}\partial_r$. 显然, 

\[
\sum_{i}|A_{i}\phi| \lesssim |\nabla\phi| \lesssim \sum_{i}|A_{i}\phi|
\]

此外, 

\[
|\partial_{r}\phi| + \sum_{i}|A_{i}\phi| \lesssim |D\phi| \lesssim |\partial_{r}\phi| + \sum_{i}|A_{i}\phi|
\]

我们还可以容易地验证如下简单的代数恒等式, 

\[
\begin{array}{rcl}
\frac{1}{2}\big((t + r)E_{+} + (t - r)E_{-}\big) & = & S \\
\frac{1}{2}\big((t + r)E_{+} - (t - r)E_{-}\big) & = & \sum_{i}\frac{x_{i}}{|x|}L_{i} \\
tA_{i} & = & L_{i} - \frac{x_{i}}{|x|}\sum_{j}\frac{x_{j}}{|x|}L_{j} \\
t\Omega_{ij} & = & x_{i}L_{j} - x_{j}L_{i}
\end{array}
\]

由前两个恒等式容易导出, 

\begin{equation}
\begin{array}{l}
|E_{+}\phi(t,  x)| \lesssim \frac{1}{t}|\Gamma\phi(t,  x)| \\
|E_{-}\phi(t,  x)| \lesssim \frac{1}{|t - |x||}|\Gamma\phi(t,  x)|
\end{array}\label{eq:cG-102}
\end{equation}

其中 $|\Gamma\phi| = |S\phi| + |L\phi|$. 

\begin{equation}
|\nabla\phi(t,  x)| \lesssim \frac{1}{t}|\Gamma\phi(t,  x)|.\label{eq:cG-103}
\end{equation}

显然, 我们还有

\[
|\partial\phi(t,  x)| \lesssim \frac{1}{|t - |x||}|\Gamma\phi(t,  x)|
\]

或更一般地, 

\begin{equation}
|\partial^{N}\phi(t,  x)| \lesssim \frac{1}{\left|t - |x|\right|^{N}}|\Gamma^{N}\phi(t,  x)|\label{eq:cG-104}
\end{equation}

其中 $|\Gamma^N\phi| = \sum |\Gamma_1\ldots \Gamma_N\phi|$, $\Gamma_1, \dots , \Gamma_N$ 是向量场 $S, L_{1}, \ldots L_{n}$ 中的任意一个. 

将上述不等式与我们范数的定义结合, 导出

\[
\begin{array}{rcl}
t\|E_{+}\phi(t)\|_{L^{2}} & \lesssim & E_{1}(\phi)^{\frac{1}{2}} \\
t\|\nabla\phi(t)\|_{L^{2}} & \lesssim & E_{1}(\phi)^{\frac{1}{2}} \\
\|uE_{-}\phi(t)\|_{L^{2}} & \lesssim & E_{1}(\phi)^{\frac{1}{2}}
\end{array}
\]

我们现在准备证明如下结论:

\begin{proposition}\label{prop:cG-4.15}
设 $\square \phi = 0$, 其初值数据满足上述假设. 则对所有 $t \geq \Lambda_0$, $s > \frac{n}{2}$, 

\begin{equation}
\|\phi(t)\|_{L^{\infty}} \lesssim \left(\frac{1}{t}\right)^{\frac{n - 1}{2}}E_{s}[\phi](t)^{\frac{1}{2}}\label{eq:cG-105}
\end{equation}

\begin{equation}
\|(1 + |u|)^{k}\partial^{k}\phi(t)\|_{L^{\infty}} \lesssim \left(\frac{1}{t}\right)^{\frac{n - 1}{2}}E_{s + k}[\phi](t)^{\frac{1}{2}}\label{eq:cG-106}
\end{equation}

又有

\[
\begin{array}{r}
\|E_{+}\phi(t)\|_{L^{\infty}} \lesssim \left(\frac{1}{t}\right)^{\frac{n + 1}{2}}E_{s + 1}[\phi](t)^{\frac{1}{2}} \\
\|\nabla\phi(t)\|_{L^{\infty}} \lesssim \left(\frac{1}{t}\right)^{\frac{n + 1}{2}}E_{s + 1}[\phi](t)^{\frac{1}{2}} \\
\|(1 + |u|)E_{-}\phi(t)\|_{L^{\infty}} \lesssim \left(\frac{1}{t}\right)^{\frac{n - 1}{2}}E_{s + 1}[\phi](t)^{\frac{1}{2}}
\end{array}
\]
\end{proposition}

证明基于如下引理:

\begin{lemma}\label{lem:cG-4.16}
设 $u(x)$ 是 $\mathbb{R}^n$($n \geq 2$)上的光滑、紧支集函数. 我们有

\begin{equation}
|u(x)| \leq C\frac{1}{|x|^{n - 1}}\left(\|\partial_{r}u\|_{L^{1}} + \|(r\nabla)^{n - 1}\partial_{r}u\|_{L^{1}}\right)\label{eq:cG-107}
\end{equation}
\end{lemma}

作为该引理的直接推论, 我们推得

\begin{equation}
|u(x)| \leq C\frac{1}{|x|^{n - 1}}\sum_{k = 0}^{n - 1}\|\partial_{r}\Omega^{k}u\|_{L^{1}}\label{eq:cG-108}
\end{equation}

其中 $\Omega^{k}u$ 表示所有形如 $\Omega_{1}\ldots\Omega_{k}u$ 的导数, $\Omega_{i}$ 是角动量向量场 $\Omega_{i}=\epsilon_{ijk}x_{j}\partial_{k}$. 现在把不等式 \eqref{eq:cG-108} 应用于函数 $u(x)=\phi^{2}(t, x)$, 其中某个 $t$ 固定. 

\[
\begin{array}{rcl}
|\phi(t,  x)|^{2} & \lesssim & \frac{1}{|x|^{n - 1}}\sum_{k = 0}^{n - 1}\sum_{i + j = k}\|\Omega^{i}\phi(t)\|_{L^{2}}\|\partial\Omega^{j}\phi(t)\|_{L^{2}} \\
& \lesssim & \frac{1}{|x|^{n - 1}}E_{n}[\phi]
\end{array}
\]

由此我们推知, 对 $|x| \geq \frac{t}{2}$, 

\begin{equation}
|\phi(t,  x)| \lesssim \left(\frac{1}{t}\right)^{\frac{n - 1}{2}}E_{n}^{\frac{1}{2}}[\phi]\label{eq:cG-109}
\end{equation}

现在考虑区域 $|x| \leq \frac{t}{2}$. 设 $\chi$ 是支集在区域 $|x| \leq 1$ 内的检验函数, 满足 $\chi(x) = 1$(若 $|x| \leq \frac{1}{2}$). 把 Sobolev 不等式

\[
|u(x)| \lesssim \|D^{n}u\|_{L^{1}}.
\]

应用于函数 $u(x) = \chi(x/t)\phi^{2}(t,  x)$. 于是, 对 $|x| \leq \frac{t}{2}$, 并考虑到 \eqref{eq:cG-104} 与 \eqref{eq:cG-31}, 

\[
\begin{array}{rcl}
|\phi(t,  x)|^{2} & \lesssim & \sum_{i + j + k = n}\|D^{i}\chi(x/t)\|_{L_{\infty}}\cdot\|D^{j}\phi(t)\|_{L^{2}(|x| \leq \frac{t}{2})}\cdot\|D^{k}\phi(t)\|_{L^{2}(|x| \leq \frac{t}{2})} \\
& \lesssim & \sum_{i + j + k = n}t^{-i}\cdot t^{-j}\|\Gamma^{j}\phi(t)\|_{L^{2}}\cdot t^{-k}\|\Gamma^{k}\phi(t)\|_{L^{2}} \\
& \lesssim & t^{-n}E_{n}[\phi]
\end{array}
\]

因此, 

\[
\sup_{|x| \leq \frac{t}{2}}|\phi(t,  x)| \lesssim \left(\frac{1}{t}\right)^{\frac{n}{2}}E_{n}^{\frac{1}{2}}[\phi].
\]

它与 \eqref{eq:cG-109} 一起证明了

\[
\|\phi(t)\|_{L^{\infty}} \lesssim \left(\frac{1}{t}\right)^{\frac{n - 1}{2}}E^{\frac{1}{2}}[\phi]
\]

这正是所期望的. 

为证明 \lemref{lem:cG-4.16}, 在极坐标 $x = r\xi$($\xi \in \mathbb{R}^{n - 1}$)下写出

\[
u(r\xi) = -\int_{r}^{\infty}\partial_{r}u(\lambda\xi)d\lambda
\]

于是

\[
\int_{|\xi| = 1}|u(r\xi)|d\sigma(\xi) \leq \frac{1}{r^{n - 1}}\int_{\mathbb{R}^{n}}|\partial_{r}u(y)|dy.
\]

另一方面, 利用单位球面 $S^{n - 1}$ 上的 Sobolev 不等式, 

\[
|u(x)| \leq c_{n}\left(\|u(r\cdot)\|_{L^{1}\left(S^{n - 1}\right)} + \|(r\nabla)^{n - 1}u(r\cdot)\|_{L^{1}\left(S^{n - 1}\right)}\right)
\]

它与上述不等式结合即证明所期望的结果. 

\begin{definition}\label{def:cG-4.17}
设 $\phi,  \psi$ 是 $\mathbb{R}^{n+1}$ 上的两个光滑函数. 我们定义

\[
\begin{array}{rcl}
Q_{0}(\phi,  \psi) & = & m^{\alpha\beta}\partial_{\alpha}\phi\, \partial_{\beta}\psi \\
Q_{\alpha\beta}(\phi,  \psi) & = & \partial_{\alpha}\phi\, \partial_{\beta}\psi - \partial_{\alpha}\psi\, \partial_{\beta}\phi
\end{array}
\]
\end{definition}

下面我们将利用之前建立的衰减估计以及零形式在 {[}28{]} 精神下的适当分解, 来估计零二次型 $Q_{0}(\phi, \psi)$ 与 $Q_{ij}(\phi, \psi)$. 假设 $\Box\phi=\Box\psi=0$. 考虑零形式 $Q_{0}(\phi, \psi)$, 将其写成

\[
Q_{0}(\phi,  \psi) = -\frac{1}{2}(E_{+}\phi\cdot E_{-}\psi + E_{-}\phi\cdot E_{+}\psi) + \nabla\phi\cdot\nabla\psi
\]

现在, 根据 \propref{prop:cG-4.15} 的结果, 

\[
\begin{array}{rcl}
\|Q_{0}(\phi,  \psi)(t)\|_{L^{2}} & \leq & \frac{1}{2}\bigg(\|E_{+}(\phi)(t)\|_{L^{\infty}}\|E_{-}(\psi)(t)\|_{L^{2}} \\
& & + \|E_{-}(\phi)(t)\|_{L^{\infty}}\|E_{+}(\psi)(t)\|_{L^{2}}\bigg) + \|\nabla\phi(t)\|_{L^{\infty}}\|\nabla\psi(t)\|_{L^{2}} \\
& \lesssim & \left(\frac{1}{t}\right)^{\frac{n + 1}{2}}E_{n + 1}[\phi]^{\frac{1}{2}}E_{1}[\psi]^{\frac{1}{2}}
\end{array}
\]

类似地, 为估计 $Q_{ij}(\phi , \psi)$, 我们写出 $\partial_i = A_i + \frac{x_i}{|x|}\partial_r$. 因此, 

\[
Q_{ij}(\phi,  \psi) = A_{i}\phi A_{j}\psi - A_{j}\phi A_{i}\psi + \partial_{r}\phi\Big(\frac{x_{i}}{r}A_{j}\phi - \frac{x_{j}}{r}A_{i}\phi\Big)
\]

于是, 考虑到 \eqref{eq:cG-32} 中的第三个公式以及 \eqref{eq:cG-31}, 我们推出 $Q_{ij}$ 具有同类型的估计. 
\subsection{第二讲的补充}

\mtsubsecn{5.0.1}{波映射}

从 Minkowski 时空到 Riemann 流形 $(M, g)$ 的波映射是映射 $\phi:\mathbb{R}^{1+n}\to M$, 它是 Lagrange 量

\[
\mathcal{L}[u] = \frac{1}{2}\int_{\mathbb{R}^{1+n}}\langle du,  du\rangle\, dt\, dx, 
\]

关于紧支集变分的临界点, 其中 $\langle du,  du\rangle = \sum_{\mu=0}^{n}\sum_{a, b} g_{ab}\, \partial_{\mu}u^{a}\partial^{\mu}u^{b}$, 这里用的是 $M$ 上的局部坐标. 该变分问题的 Euler-Lagrange 方程, 在 $M$ 的局部坐标下恰好具有 \eqref{eq:cG-WM} 的形式, 其中 $\Gamma_{JK}^{I}$ 是局部坐标下 $M$ 的 Christoffel 符号, $N = \dim M$(见, 例如, Shatah-Struwe {[}58{]}). 

\mtsubsecn{5.0.2}{Maxwell-Klein-Gordon 方程组}

在下面的讨论中, 求和约定生效. 希腊指标从 0 到 $n$ 求和, 罗马指标从 1 到 n. 回顾指标相对于 Minkowski 度量 $m_{\mu\nu} = \text{diag}(-1,  1,  \ldots,  1)$ 升降. 例如, $\Box = \partial^{\mu}\partial_{\mu}$, $\Delta = \partial^{j}\partial_{j}$. 我们用 $i$ 表示虚数单位. 

方程组的未知量是一个一次形式 $A_{\mu}dx^{\mu}$(规范势)和一个标量 $\phi$, 二者都定义在 Minkowski 时空上:

\[
\begin{array}{c}
A_{\mu}: \mathbb{R}^{1+n} \to \mathbb{R},  \\
\phi: \mathbb{R}^{1+n} \to \mathbf{C}.
\end{array}
\]

电磁场是二次形式 $F_{\mu\nu} = \partial_{\mu}A_{\nu} - \partial_{\nu}A_{\mu}$. 相对于规范势的协变导数是

\[
D_{\mu}\phi = \partial_{\mu}\phi + iA_{\mu}\phi.
\]

我们寻找 Lagrange 量

\[
\mathcal{L}\left[A_{\mu},  \phi\right] = \int_{\mathbb{R}^{1+n}}\left(-\frac{1}{4}F_{\mu\nu}F^{\mu\nu} - \frac{1}{2}D_{\mu}\phi\, \overline{D^{\mu}\phi}\right)dt\, dx.
\]

的临界点. 相应的 Euler-Lagrange 方程是

\begin{equation}
\partial^{\mu}F_{\mu\nu} = -\Im\left(\phi\, \overline{D_{\nu}\phi}\right), \tag{MKGa}\label{eq:cG-MKGa}
\end{equation}

\begin{equation}
D^{\mu}D_{\mu}\phi = 0, \tag{MKGb}\label{eq:cG-MKGb}
\end{equation}

其中 $\Im z$ 表示 $z$ 的虚部. 

设 $\chi$ 是 $\mathbb{R}^{1+n}$ 上的实值函数, 考虑由

\[
\begin{array}{c}
\widetilde{A}_{\mu} = A_{\mu} - \partial_{\mu}\chi,  \\
\widetilde{\phi} = e^{i\chi}\phi.
\end{array}
\]

给出的变换 $(A_{\mu}, \phi)\to (\widetilde{A}_{\mu}, \widetilde{\phi})$. 显然, 电磁场在规范变换 $A_{\mu} \rightarrow \widetilde{A}_{\mu}$ 下保持不变, 而一个简单的计算表明, 若 $(A_{\mu},  \phi)$ 满足 (MKG), 则 $(\widetilde{A}_{\mu},  \widetilde{\phi})$ 也满足(记住 $D_{\mu}$ 依赖于 $A_{\mu}$). 这就在满足 (MKG) 的二元组 $(A_{\mu},  \phi)$ 的集合上给出了一个等价关系, 而所谓后者的一个解, 我们指的是这样一类二元组的一个等价类. 

于是, 我们拥有规范自由度, 即我们可以自由选择给定解(等价类)的任意代表元, 并且可以规定规范势应满足的条件. 传统的规范条件是:

\begin{itemize}
\item Lorentz 规范:$\partial^{\mu}A_{\mu}=0, $
\item Coulomb 规范:$\partial^{j}A_{j}=0$, 
\item 时间规范:$A_0 = 0$. 
\end{itemize}

\noindent\textbf{(MKG) 在 Lorentz 规范中. }把 Lorentz 条件与 (MKG) 耦合, 得到方程组

\begin{subequations}\label{eq:cG-110}
\begin{align}
\Box A_{\mu} &= -\Im\big(\phi\, \overline{\partial_{\mu}\phi}\big) + |\phi|^{2}A_{\mu},  \label{eq:cG-110a} \\
\Box \phi &= -2iA^{\mu}\partial_{\mu}\phi + A^{\mu}A_{\mu}\phi,  \label{eq:cG-110b} \\
\partial^{\mu}A_{\mu} &= 0. \label{eq:cG-110c}
\end{align}
\end{subequations}

现在注意, 若 $(A_{\mu},  \phi)$ 满足 \eqref{eq:cG-110a} 与 \eqref{eq:cG-110b}, 其初值数据为

\begin{subequations}\label{eq:cG-111}
\begin{align}
A_{\mu}|_{t = 0} = a_{\mu},  \qquad &\partial_{t}A_{\mu}|_{t = 0} = b_{\mu},  \label{eq:cG-111a} \\
\phi|_{t = 0} = \phi_{0},  \qquad &\partial_{t}\phi|_{t = 0} = \phi_{1},  \label{eq:cG-111b}
\end{align}
\end{subequations}

且满足约束条件

\begin{equation}
b_{0} = \partial^{j}a_{j},  \qquad \Delta a_{0} - |\phi_{0}|^{2}a_{0} = \partial^{j}b_{j} - \Im(\phi_{0}\overline{\phi}_{1}), \label{eq:cG-112}
\end{equation}

则 \eqref{eq:cG-110c} 自动满足. 因为由 \eqref{eq:cG-110a} 与 \eqref{eq:cG-110b} 可知, $u = \partial^{\mu}A_{\mu}$ 满足

\[
\Box u = |\phi|^{2}u, 
\]

而由 \eqref{eq:cG-111a} 与 \eqref{eq:cG-112} 可知, $u|_{t=0} = \partial_t u|_{t=0} = 0$. 由解的唯一性, $u = 0$. 

于是, \eqref{eq:cG-110c} 等价于初值数据上的约束 \eqref{eq:cG-112}, 所以我们只剩 \eqref{eq:cG-110a} 与 \eqref{eq:cG-110b}. 因此, Lorentz 规范中的 (MKG) 示意性地具有 $\Box u = u\partial u + u^{3}$ 的形式. 遗憾的是, 这类方程的一般情形没有良好的局部正则性, 所以 Lorentz 规范对我们的目的不是很有用. 

\noindent\textbf{(MKG) 在 Coulomb 规范中. }把 Coulomb 条件与 (MKG) 耦合, 给出

\begin{subequations}\label{eq:cG-113}
\begin{align}
\Delta A_{0} &= -\Im\big(\phi\, \overline{\partial_{t}\phi}\big) + |\phi|^{2}A_{0},  \label{eq:cG-113a} \\
\Box A_{j} &= -\Im\big(\phi\, \overline{\partial_{j}\phi}\big) + |\phi|^{2}A_{j} - \partial_{j}\partial_{t}A_{0},  \label{eq:cG-113b} \\
\Box \phi &= -2iA^{j}\partial_{j}\phi + 2iA_{0}\partial_{t}\phi + i(\partial_{t}A_{0})\phi + A^{\mu}A_{\mu}\phi,  \label{eq:cG-113c} \\
\partial^{j}A_{j} &= 0. \label{eq:cG-113d}
\end{align}
\end{subequations}

这里我们把规范势分解为它的时间分量 $A_{0}$ 与空间分量 $A = A_{j}dx^{j}$. 我们在时刻 t = 0 给定初值数据:

\begin{subequations}\label{eq:cG-114}
\begin{align}
A_{j}|_{t = 0} = a_{j},  \qquad &\partial_{t}A_{j}|_{t = 0} = b_{j},  \label{eq:cG-114a} \\
\phi|_{t = 0} = \phi_{0},  \qquad &\partial_{t}\phi|_{t = 0} = \phi_{1}. \label{eq:cG-114b}
\end{align}
\end{subequations}

对 $A_0$ 不施加初值条件;若我们令 $a_0 = A_0|_{t=0}$, 则由 \eqref{eq:cG-113a} 得 $\Delta a_0 - |\phi_0|^2 a_0 = -\Im(\phi_0\overline{\phi}_1)$. 

方程 \eqref{eq:cG-113d} 在数据无散时自动满足:

\begin{equation}
\partial^{j}a_{j} = \partial^{j}b_{j} = 0.\label{eq:cG-115}
\end{equation}

因为若 $(A_0,  A,  \phi)$ 满足 \eqref{eq:cG-113a}--\eqref{eq:cG-113c}, 则 $u = \partial^j A_j$ 满足 $\Box u = |\phi|^2 u$, 而若 \eqref{eq:cG-114} 与 \eqref{eq:cG-115} 成立, 则 $u|_{t=0} = \partial_t u|_{t=0} = 0$. 

于是我们只剩下方程 \eqref{eq:cG-113a}--\eqref{eq:cG-113c}. 其中第一个是椭圆方程, 相对容易处理, 所以我们把它排除在我们的模型方程之外. 剩下的两个方程右端有三类项:

\begin{itemize}
\item 涉及 $A_0$ 的"椭圆项";这些项统称为 $\mathcal{E}$. 
\item $A_{j}$ 与 $\phi$ 的三次项;这些项统称为 $\mathcal{C}$. 
\item 具有零形式结构的二次项. 
\end{itemize}

属于后一类的项是 $-\Im\big(\phi\, \overline{\partial_{j}\phi}\big)$ 与 $-2iA^{j}\partial_{j}\phi$. 我们现在揭示这些表达式(由于 Coulomb 条件)内在的零形式结构. 

把 $\phi$ 分解为实部与虚部:$\phi = u + iv$. 则

\[
-\Im\bigl(\phi\, \overline{\partial_{j}\phi}\bigr) = u\partial_{j}v - v\partial_{j}u, 
\]

所以 \eqref{eq:cG-113b} 读作 $\mathbb{R}^n$ 上(依赖于时间的)一次形式的方程:

\[
\Box A = u\, dv - v\, du + \mathcal{C} - d(\partial_{t}A_{0}).
\]

对两边作用 $d$:

\[
\Box(dA) = 2\, du \wedge dv + d\mathcal{C}.
\]

但

\[
du \wedge dv = \frac{1}{2}Q_{jk}(u,  v)\, dx^{j} \wedge dx^{k}, 
\]

由此

\[
\Box F_{jk} = Q_{jk}(u,  v) + \partial\mathcal{C}.
\]

Coulomb 规范条件意味着 $\partial^k F_{jk} = -\Delta A_j$, 所以我们有

\[
-\Delta\square A_{j} = \partial^{k}Q_{jk}(u,  v) + \partial^{2}\mathcal{C}.
\]

于是, 模去 Riesz 算子, 

\begin{equation}
\Box A = D^{-1}Q\left(\Re\phi,  \Im\phi\right) + \mathcal{C}\label{eq:cG-116}
\end{equation}

其中 $Q$ 是零形式\footnote{确切地说, $Q$ 的第 $j$ 个分量是 $\sum_{k}R_kQ_{jk}$, 其中 $R_k = D^{-1}\partial_k$ 是第 $k$ 个 Riesz 算子. 由于我们使用的范数只依赖于 傅里叶变换的大小, 我们忽略 Riesz 算子. }$Q_{jk}$ 的某个线性组合. 由于三次项 $\mathcal{C}$ 更容易估计, 我们把它排除在我们的模型问题之外. 

现在考虑方程 \eqref{eq:cG-113c}. 分离实部与虚部, 我们有

\[
\begin{array}{rcl}
\square u &=& 2A\cdot\nabla v + \mathcal{C} + \mathcal{E},  \\
\square v &=& -2A\cdot\nabla u + \mathcal{C} + \mathcal{E}.
\end{array}
\]

(这里我们通过提升指标把 $A$ 看作向量场;$\nabla$ 表示关于空间变量的梯度. )我们断言, 项 $A\cdot\nabla u$ 与 $A\cdot\nabla v$ 具有零形式结构, 其原因在于 $A$ 无散(由 Coulomb 条件). 设 $B_{jk}$ 是

\begin{equation}
\Delta B_{jk} = \partial_{j}A_{k} - \partial_{k}A_{j}\label{eq:cG-117}
\end{equation}

的唯一解(带适当的正则性假设). 由 Coulomb 条件, 

\begin{equation}
\Delta\partial^{j}B_{jk} = \Delta A_{k},  \quad \text{这意味着} \quad \partial^{j}B_{jk} = A_{k}\label{eq:cG-118}
\end{equation}

于是, 

\[
A\cdot\nabla u = \partial^{j}B_{jk}\partial^{k}u = \frac{1}{2}Q_{jk}(u,  B^{jk}).
\]

因此, $u = \Re\phi$ 与 $v = \Im\phi$ 的上述方程可以改写为

\[
\begin{array}{rcl}
\Box\Re\phi &=& Q_{jk}(\Im\phi,  B^{jk}) + \mathcal{C} + \mathcal{E},  \\
\Box\Im\phi &=& Q_{jk}(B^{jk},  \Re\phi) + \mathcal{C} + \mathcal{E}.
\end{array}
\]

但根据 \eqref{eq:cG-117}, 模去 Riesz 算子, $B$ 具有 $D^{-1}A$ 的形式. 将此与 \eqref{eq:cG-116} 结合, 并处处丢弃项 $\mathcal{C}$ 与 $\mathcal{E}$, 我们得到形如 ("MKG") 的方程组, 它就是 (MKG) 的模型. 

\mtsubsecn{5.0.3}{Yang-Mills 方程组}

设 $G$ 是经典紧致矩阵 李群之一(如 $SO(k, \mathbb{R})$ 或 $SU(k, \mathbf{C})$), $\mathfrak{g}$ 是其 李代数. 未知量是 $R^{1+n}$ 上的一个 $\mathfrak{g}$ 值一次形式 $A_{\mu}dx^{\mu}$. 相应的协变导数是

\[
D_{\mu}H = \partial_{\mu}H + [A_{\mu},  H], 
\]

其中 $H$ 是 $\mathbb{R}^{1+n}$ 上的任意 $\mathfrak{g}$ 值张量场, $[\cdot, \cdot]$ 是矩阵换位子. 

曲率是 $\mathfrak{g}$ 值二次形式

\[
F_{\mu\nu} = \partial_{\mu}A_{\nu} - \partial_{\nu}A_{\mu} + [A_{\mu},  A_{\nu}].
\]

Lagrange 量是

\[
\mathcal{L}[A_{\mu}] = -\frac{1}{4}\int\langle F_{\mu\nu},  F^{\mu\nu}\rangle\, dt\, dx, 
\]

其中 $\langle\cdot, \cdot\rangle$ 是 $\mathfrak{g}$ 从环境空间继承的内积(例如, $SO(k, \mathbb{R})$ 嵌入 $\mathbb{R}^{k^2}$, 所以它的 李代数可以看作后者的子空间). Euler-Lagrange 方程是

\[
D^{\nu}F_{\mu\nu} = 0.\tag{YM}
\]

设 $O$ 是 $\mathbb{R}^{1+n}$ 上的 $G$ 值函数. 考虑由

\[
\widetilde{A}_{\mu} = OA_{\mu}O^{-1} - \partial_{\mu}O\, O^{-1}.
\]

给出的规范变换 $A_{\mu} \to \widetilde{A}_{\mu}$. 一个计算表明, 曲率于是变换为

\[
\widetilde{F}_{\mu\nu} = OF_{\mu\nu}O^{-1}.
\]

记 $\widetilde{D}_{\mu}$ 为对应于 $\widetilde{A}_{\mu}$ 的协变导数, 则我们有

\[
\widetilde{D}^{\nu}\widetilde{F}_{\mu\nu} = OD^{\nu}F_{\mu\nu}O^{-1}, 
\]

所以 (119) 在规范变换下不变. 因此我们拥有规范自由度, 可以对 $A_{\mu}$ 施加规范条件. 

\noindent\textbf{(YM) 在 Coulomb 规范中. }相对于 Coulomb 条件 $\partial^j A_j = 0$, (119) 采取如下形式(见 {[}34{]})

\begin{subequations}\label{eq:cG-119}
\begin{align}
\Delta A_{0} &= 2[\partial^{j}A_{0},  A_{j}] + [A^{j},  \partial_{t}A_{j}] + [A^{j},  [A_{0},  A_{j}]],  \label{eq:cG-119a} \\
\Box A_{j} + \partial_{t}\partial_{j}A_{0} &= -2[A^{k},  \partial_{k}A_{j}] + [A^{k},  \partial_{j}A_{k}] + [\partial_{t}A_{0},  A_{j}] + 2[A_{0},  \partial_{t}A_{j}] \nonumber \\
&\quad - [A_{0},  \partial_{j}A_{0}] - [A^{k},  [A_{k},  A_{j}]] + [A_{0},  [A_{0},  A_{j}]],  \label{eq:cG-119b} \\
\partial^{j}A_{j} &= 0. \label{eq:cG-119c}
\end{align}
\end{subequations}

遗憾的是, 假设全局 Coulomb 规范的存在迫使对初值数据施加一个限制性很强的小性假设. 在 {[}34{]} 中, 这个困难通过使用局部论证得以解决. 沿用 {[}47{]}, 我们忽略这一复杂性, 并从方程组 (119) 导出我们的模型方程. 

与 (MKG) 的讨论一样, (119c) 归结为初值数据上的一个约束. $A_{0}$ 的方程是椭圆的, 所以我们忽略它. 至于 (119b), 我们只保留右端前两项, 因为其余各项要么涉及 $A_{0}$(对 $A_0$ 我们期望有比 $A_j$ 更好的估计), 要么是三次项. 

现在把 (119b) 写成 $\mathbb{R}^n$ 上依赖于时间的 $\mathfrak{g}$ 值一次形式的方程(忽略右端除前两项外的所有项):

\[
\Box A + d(\partial_{t}A_{0}) = S + T, 
\]

其中 $A = A_{j}dx^{j}$, $S = -2[A^{k}, \partial_{k}A_{j}]dx^{j}$, $T = [A^k, \partial_jA_k]dx^j$. 对两边作用外微分 $d$:

\[
\square dA = dS + dT.
\]

设 $B$ 是由方程 \eqref{eq:cG-117} 确定的二次形式(此时是 $\mathfrak{g}$ 值的). 于是

\[
\Box\Delta B_{jk} = \partial_{j}S_{k} - \partial_{k}S_{j} + \partial_{j}T_{k} - \partial_{k}T_{j}.
\]

由 \eqref{eq:cG-118}, 推出

\[
-\Box\Delta A_{j} = \partial^{k}\left(\partial_{j}S_{k} - \partial_{k}S_{j} + \partial_{j}T_{k} - \partial_{k}T_{j}\right), 
\]

所以, 出于频率空间估计的目的, 我们可以用

\begin{equation}
\square A = S + D^{-1}dT\label{eq:cG-120}
\end{equation}

替换 (119b). 

剩下的工作是识别隐藏在 $S$ 与 $dT$ 中的零形式结构. 首先, 我们有

\[
\begin{array}{rcl}
S_{j} &=& [\partial_{k}B^{kl},  \partial_{l}A_{j}] = \frac{1}{2}[\partial_{k}B^{kl} - \partial_{k}B^{lk},  \partial_{l}A_{j}] \\
&=& \frac{1}{2}\left([\partial_{k}B^{kl},  \partial_{l}A_{j}] - [\partial_{l}B^{kl},  \partial_{k}A_{j}]\right) \\
&=& \frac{1}{2}\left(\partial_{k}B^{kl}\partial_{l}A_{j} - \partial_{l}B^{kl}\partial_{k}A_{j} + \partial_{k}A_{j}\partial_{l}B^{kl} - \partial_{l}A_{j}\partial_{k}B^{kl}\right), 
\end{array}
\]

所以矩阵 $S_{j}$ 的每个元素都是形如 $Q_{kl}(A,  B)$ 的项的线性组合, 其中 $A$ 与 $B$ 代表 $A_{j}$ 与 $B_{kl}$ 的任意两个元素. 但由 \eqref{eq:cG-118}, 我们可以用 $D^{-1}A$ 替换 B. 示意性地, 

\begin{equation}
S \sim Q(A,  D^{-1}A)\label{eq:cG-121}
\end{equation}

现在考虑一次形式 $T$. 我们计算:

\[
\begin{array}{rcl}
(dT)_{jk} &=& \partial_{j}[A^{l},  \partial_{k}A_{l}] - \partial_{k}[A^{l},  \partial_{j}A_{l}] \\
&=& [\partial_{j}A^{l},  \partial_{k}A_{l}] - [\partial_{k}A^{l},  \partial_{j}A_{l}].
\end{array}
\]

于是, 矩阵 $(dT)_{jk}$ 的每个元素都是形如 $Q_{jk}(A, A')$ 的项的线性组合, 其中 $A$ 与 $A'$ 代表 $A_{l}$($1 \leq l \leq n$)的任意两个元素. 将此与 \eqref{eq:cG-121} 和 \eqref{eq:cG-120} 结合, 我们得到 Yang-Mills 方程组的模型 ("YM"). 
\subsection{第三讲的补充}

该结果的一个重要推论如下\footnote{这一点最近在 {[}40{]} 中给出了直接证明. }:

\begin{theorem}\label{thm:cG-6.1}
对任意渐近平坦初值数据集 $(\Sigma, \bar{g}, \bar{k})$, 设 $(\mathcal{M}, g)$ 是其最大未来发展, 则可以在 $\Sigma$ 中找到紧闭包 (compact closure) 的适当区域 $\Omega_{0}$, 使得其在 $\mathcal{M}$ 中的影响域\footnote{也称为 $\Omega_{0}$ 的因果未来}$\mathcal{C}^{+}(\Omega_{0})$ 的边界 $D_{0}^{+}$ 具有完整的零测地生成线\footnote{相对于相应的仿射参数}. 
\end{theorem}

上述推论可以用来引入完整未来零无穷大的概念\footnote{这个概念在广义相对论文献中通常通过正则共形紧化的概念定义, 即在无穷远处粘贴一个边界. 这里给出的定义来自 {[}10{]}, 它避免了紧化精确光滑程度这一技术问题. }:

\begin{definition}\label{def:cG-6.2}
我们说 AF 初值数据集 $(\Sigma,  \bar{g},  \bar{k})$ 的最大未来发展 $(\mathcal{M},  g)$ 具有完整的未来零无穷大, 若对任意正实数 A, 都可以找到包含定理 1.2 的集合 $\Omega_{0}$ 的区域 $\Omega$, 使得 $\Omega$ 在 $\mathcal{M}$ 中的依赖域的边界 $D^{-}$ 具有如下性质:它的每条零测地生成线, 从 $D^{-} \cap D_{0}^{+}$ 量起, 总仿射长度至少为 A. 
\end{definition}

对足够大的初值数据集而言, 奇点的不可避免的存在性促使 Penrose 提出了两个猜想, 即所谓的弱宇宙监督猜想与强宇宙监督猜想. 

\noindent\textbf{猜想 1}(弱宇宙监督猜想(WCC)). 一般的渐近平坦初值数据具有拥有完整未来零无穷大的最大未来发展. 

WCC 猜想并不排除奇点可能被局部观测者看见的可能性. 这可能导致这样的悖论情形:这些观测者所作的观测结果缺乏唯一可预测性. 强宇宙监督正是为禁止奇点的这种不良特征而设计的. 

\noindent\textbf{猜想 2}(强宇宙监督). 一般的初值数据\footnote{不一定渐近平坦. 这一猜想常在紧致初值数据集的背景下讨论. }集具有作为 Lorentz 流形在连续意义下局部不可延拓的最大未来发展. 

用更技术性的语言来说, 这意味着, 撇开某些可能的例外初值条件不论, 初值数据集的最大未来发展具有如下性质:沿着任意未来方向、不可延拓、有限长度\footnote{即有界固有时间}的类时测地线, 时空曲率分量(相对于沿测地线平行移动的标准正交标架表示)必定在弧长趋于其极限值时变为无穷大. 

\mtsubsec{\thmref{thm:cG-3.10} 到\thmref{thm:cG-3.12} 的化简}

下面我们描述 \thmref{thm:cG-3.10} 的化简, 从现在起我们把它称为定理 (A);化简的目标是证明 \thmref{thm:cG-3.12} 的衰减估计, 后者将称为定理 (B). 前五步现在是标准的, 见 {[}2{]}, 尤其见 {[}72{]}. 到最后一步的化简是 {[}29{]} 的几何方法的典型特征. 

\mtsubsec{第 1 步：能量估计}

正如我们在上面已经注意到的, 拟线性波动方程的能量估计意味着, 解 $\phi$ 的 $L_{[0, T]}^{\infty}H^{s}$ 范数受初值数据 $\phi[0]$ 的 $H^{s}$ 范数控制, 只要 $\|\partial\phi\|_{L_{[0, T]}^{1}L_{x}^{\infty}}\leq 1$. 下一个命题是该陈述的一个更精确的版本. 

\begin{proposition}[能量估计]\label{prop:cG-6.5}
设 $\phi \in C([0, T], H^s)\cap C^1([0, T], H^{s - 1})$ 是 \eqref{eq:cG-67} 在时间区间 $[0, T]$ 上对某个 $s\geq 1$ 的解, 满足条件 $\|\phi\|_{L_{[0, T]}^{\infty}L_x^\infty}\leq \Lambda_0$. 则 $\phi$ 满足如下能量估计. 

\begin{equation}
\left\|\phi\right\|_{L_{[0, T]}^{\infty}\dot{H}^{s}} \leq C\left(\left\|\partial\phi\right\|_{L_{[0, T]}^{1}L_{x}^{\infty}},  \Lambda_{0}\right)\|\phi[0]\|_{\dot{H}^{s}}.\label{eq:cG-122}
\end{equation}

这里常数 $C(a, b)$ 表示对 $a, b$ 的依赖, 而 $\dot{H}^s$ 是通常的齐次 Sobolev 空间. 
\end{proposition}

\mtsubsec{第 2 步：化归为 Strichartz 型估计}

由 Cauchy-Schwarz 不等式, $\|\partial\phi\|_{L_{[0, T]}^1 L_x^\infty} \leq T^{\frac{1}{2}} \|\partial\phi\|_{L_{[0, T]}^2 L_x^\infty}$. 于是, Strichartz 不等式 \eqref{eq:cG-71} 与时间区间 $[0, T]$ 的小性可用于在空间 $H^{s*}$ 中封闭能量估计. 这实际上给出了问题所期望的局部适定性. 于是定理 (A) 可以化归为如下自举论证. 

\begin{theorem}[A1]\label{thm:cG-6.7}
设 $\phi \in C([0, T], H^{s*})\cap C^1([0, T], H^{s*-1})$ 是 \eqref{eq:cG-67} 在时间区间 $[0, T], T\leq 1$ 上的解. 假设

\begin{equation}
\left\|\phi\right\|_{L_{[0, T]}^{\infty}H^{s *}} + \left\|\partial\phi\right\|_{L_{[0, T]}^{2}L_{x}^{\infty}} \leq B_{0}, \label{eq:cG-123}
\end{equation}

其中常数 $B_0 \leq c_{s_*}^{-1}\Lambda_0$, $c_{s_*}$ 是嵌入 $H^{s_*} \subset L^\infty$ 的 Sobolev 常数. 则 $\phi$ 满足如下局部时间的 Strichartz 型估计, 

\begin{equation}
\left\|\partial\phi\right\|_{L_{[0, T]}^{2}L_{x}^{\infty}} \leq C(B_{0})T^{s_{*} - s_{0}}\left\|\phi\right\|_{L_{[0, T]}^{\infty}H^{s_{*}}}.\label{eq:cG-124}
\end{equation}
\end{theorem}

\mtsubsec{第 3 步：Strichartz 型估计的二进版本与仿微分逼近}

为了证明 Strichartz 型估计 \eqref{eq:cG-124}, 我们可以把拟线性问题 \eqref{eq:cG-67} 看作关于 $\phi$ 的具有粗糙系数的线性波动方程. 光滑化系数并处理依赖于参数的光滑系数线性波动方程族是有利的. 首先我们引入函数 $\phi^{\lambda}$, 它由 $\phi$ 限制到 Fourier 空间中频率 $\sim \lambda$ 的二进片得到. 更精确地说, 设 $\zeta$ 是支集在壳层 $\{\xi : \frac{1}{2} \leq |\xi| \leq 2\}$ 内的光滑函数. 这里, $\xi$ 表示空间 傅里叶变换的变量. 设 $\zeta$ 还满足条件 $\sum_{k \in \mathbb{Z}} \zeta(2^k\xi) = 1$, $\forall \xi \in \mathbb{R}^{3} \setminus \{\mathcal{V}\}$. 设 $\lambda$ 是二进参数 $\lambda = 2^k$(其中某个 $k \in \mathbb{Z}$), 记 $P_{\lambda}$ 为"投影算子"

\[
P_{\lambda}f(x) = f^{\lambda}(x) = \int e^{-ix\cdot\xi}\zeta(\lambda^{-1}\xi)\hat{f}(\xi)d\xi.
\]

还定义

\[
f_{\leq \lambda} = S_{\lambda}f = \sum_{\mu \leq \lambda}P_{\mu}f.
\]

于是定理 (A1) 由如下关于 $\phi^{\lambda} = P_{\lambda}\phi$ 的 Strichartz 型估计的二进版本推出. 

\begin{theorem}[A2]\label{thm:cG-6.9}
设 $\phi$ 如定理 (A1) 中所述. 固定一个大参数 $\Lambda$. 则对每个 $\lambda \geq \Lambda$, 函数 $\phi^{\lambda}$ 满足 Strichartz 型估计

\begin{equation}
\|\partial\phi^{\lambda}\|_{L_{[0, T]}^{2}L_{x}^{\infty}} \leq C(B_{0})c_{\lambda}T^{s_{*} - s_{0}}\|\phi\|_{L_{[0, T]}^{\infty}H^{s *}}, \label{eq:cG-125}
\end{equation}

其中常数 $c_{\lambda}$ 满足 $\sum_{\lambda}c_{\lambda}^{2}\leq 1$. 类似的估计对 $\phi_{\leq \Lambda}$ 也成立. 
\end{theorem}

\begin{remark*}[A2]
在低频的情形下, 关于 $\phi_{\leq \Lambda}$ 的估计 \eqref{eq:cG-125} 平凡地由 Sobolev 不等式推出. 

\[
\|\partial\phi_{\leq \Lambda}\|_{L_{[0, T]}^{2}L_{x}^{\infty}} \leq cT^{\frac{1}{2}}\|\phi_{\leq \Lambda}\|_{L_{[0, T]}^{\infty}H^{\frac{5}{2} + \epsilon}} \leq c\Lambda^{\frac{5}{2} + \epsilon - s_{*}}T^{\frac{1}{2}}\|\phi\|_{L_{[0, T]}^{\infty}H^{s_{*}}}, 
\]

其中 $c$ 是嵌入 $H^{\frac{3}{2} +\epsilon}(\mathbb{R}^3)\subset L^{\infty}(\mathbb{R}^{3})$ 的范数. 由于假设 $s_*$ 足够接近 $s_0 = 2 + \frac{1}{2}(2 - \sqrt{3})$, 且 $\Lambda$ 是只可能依赖于 $B_0$ 的固定大参数, 我们得到 $\phi$ 低频部分所期望的界. 
\end{remark*}

我们把注意力限制在大频率 $\lambda \geq \Lambda$ 上. 在下一个命题中我们证明, 对参数 $a \in [0, 1]$ 的任意固定值, 每个 $\phi^{\lambda}$ 都满足一个以光滑度量 $g_{\lambda^{a}}^{ij} = S_{\lambda^{a}}g^{ij}(S_{\lambda^{a}}\phi)$ 为非齐次波动方程的方程. 

\begin{proposition}\label{prop:cG-6.10}
设 $\phi$ 如定理 (A1) 中所述. 固定参数 $a$ 的值, $a \in [0, 1]$. 则对每个 $\lambda \geq \Lambda$, $\phi^{\lambda}$ 满足方程

\begin{equation}
\begin{aligned}
\Box_{g_{\leq\lambda^{a}}}\phi^{\lambda} &= -\partial_{t}^{2}\phi^{\lambda} + g_{\leq\lambda^{a}}^{ij}\partial_{i}\partial_{j}\phi^{\lambda} = R_{\lambda}^{a},  \\
\phi^{\lambda}|_{t = 0} &= \phi_{0}^{\lambda},  \quad \partial_{t}\phi^{\lambda}|_{t = 0} = \phi_{1}^{\lambda}.
\end{aligned}\label{eq:cG-126}
\end{equation}

此外, 右端 $R_{\lambda}^{a}$ 的 Fourier 支集包含在集合 $\{\xi : \lambda \leq |\xi| \leq 4\lambda\}$ 内, 且对所有实数 $s > 1$ 与任意 $t \in [0, T]$

\begin{equation}
\|R_{\lambda}(t)\|_{\dot{H}^{s}} \leq c\lambda^{s}\|R_{\lambda}\|_{L_{x}^{2}} \lesssim C(B_{0})\lambda^{1 - a}c_{\lambda}\|\partial\phi\|_{L_{x}^{\infty}}\|\phi\|_{\dot{H}^{s}}\label{eq:cG-127}
\end{equation}

其中常数 $c_{\lambda}$ 满足 $\sum_{\lambda} c_{\lambda}^{2} \leq 1$. 
\end{proposition}

\mtsubsec{第 4 步：频率依赖区间上的 Strichartz 估计}

这一步把定理 (A) 的证明化归为线性化方程 $\Box_{g_{\lambda^a}} \psi = 0$ 在大小 $\approx T\lambda^{-(1-a)}$ 的小时间区间 $I$ 上的精确\footnote{如平坦情形那样, 无损失}Strichartz 估计的证明. 最终 Strichartz 估计中的正则性损失, 是对 $[0,  T]$ 中的区间 $I$ 求和这些锐利 Strichartz 估计的结果. 

\begin{theorem}[A3]\label{thm:cG-6.12}
设 $\lambda \geq \Lambda$, $\psi$ 是线性波动方程 $\Box_{g_{\lambda^a}}\psi = 0$ 的解, 其初值数据 $\psi [0]$ 满足 supp $\hat{\psi}[0] \in \{\frac{1}{2}\lambda \leq |\xi| \leq 2\lambda\}$, 度量 $g_{\lambda^a}$ 如 \propref{prop:cG-6.10} 中所定义. 固定参数 $a$ 的值, $a = \frac{5 - 2s_*}{1 - 2s_* + 2s_0} < -1 + \sqrt{3}$. 则存在区间 $[0,  T]$ 的一个划分 $\{I\}$, 分成子区间 $I$, 使得每个 $I$ 的大小 $|I| \leq T\lambda^{-(1 - a)}$, 子区间的数目大约为 $\lambda^{1 - a}$, 且在每个 $I$ 上, $\psi$ 对某个固定、任意小的 $\epsilon > 0$ 满足 Strichartz 估计

\begin{equation}
\left\|P_{\lambda}\partial\psi\right\|_{L_{I}^{2}L_{x}^{\infty}} \leq C(B_{0})|I|^{\epsilon}\|\psi[0]\|_{\dot{H}^{2 + \epsilon}}\label{eq:cG-128}
\end{equation}
\end{theorem}

\begin{remark*}[A3]
根据自举假设 \eqref{eq:cG-123}, 拟线性问题的解 $\phi$ 满足估计 $\|\partial\phi\|_{L^{2}_{[0, T]} L_x^\infty} \leq B_0$. 容易推出, 存在时间区间 $[0, T]$ 的一个子划分 $\{I\}$, 子区间 $I$ 的总数介于 $\lambda^{1 - a}$ 与 $2\lambda^{1 - a}$ 之间, 每个 $I$ 的大小以 $T\lambda^{-(1 - a)}$ 为界, 使得在每个 $I$ 上有

\begin{equation}
\|\partial\phi\|_{L_{I}^{2}L_{x}^{\infty}} \leq \lambda^{-\frac{1 - a}{2}}\|\partial\phi\|_{L_{[0, T]}^{2}L_{x}^{\infty}}\label{eq:cG-129}
\end{equation}

这一构造定义了定理 (A3) 中提到的子划分 $\{I\}$. 
\end{remark*}

\mtsubsec{度量 $g_{\lambda^{a}}$ 的性质}

不等式 \eqref{eq:cG-128} 是变系数(度量)波动方程 $\Box_{g_{\leq\lambda^{a}}}\psi=0$ 的解在频率依赖区间 $I$ 上的 Strichartz 估计. 

度量 $g_{\leq\lambda^{a}} = S_{\lambda^{a}}g(S_{\lambda^{a}}\phi)$ 依赖于拟线性问题的解 $\phi$. 在下一个命题中我们陈述族 $g_{\leq\lambda^{a}}$ 的性质, 它们由 $\phi$ 上的自举条件 \eqref{eq:cG-123} 与注 (A3) 中描述的划分 $\{I\}$ 的构造推出. 

\begin{proposition}\label{prop:cG-6.14}
设 $\phi \in C([0, T], H^{s_*})\cap C^1([0, T], H^{s_* - 1})$ 是 \eqref{eq:cG-67} 在时间区间 $[0, T], T\leq 1$ 上的解. 假设 $\phi$ 满足定理 (A1) 的假设 \eqref{eq:cG-123}. 考虑注 (A3) 定义的时间区间 $[0, T]$ 的子划分 $\{I\}$. 则度量族 $g_{\leq \lambda^a} = S_{\lambda^a}g(S_{\lambda^a}\phi)$ 满足如下条件:

对所有子区间 $I$ 与所有非负整数 $m$, 

\begin{equation}
\left\|\partial^{1 + m}g_{\leq\lambda^{a}}\right\|_{L_{I}^{1}L_{x}^{\infty}} \leq \lambda^{-(1 - a) + am}\bar{B}_{0}, \label{eq:cG-130}
\end{equation}

\begin{equation}
\|\partial^{1 + m}g_{\leq\lambda^{a}}\|_{L_{I}^{2}L_{x}^{\infty}} \leq \lambda^{-\frac{1 - a}{2} + am}\bar{B}_{0}, \label{eq:cG-131}
\end{equation}

\begin{equation}
\|\partial^{1 + m}g_{\leq\lambda^{a}}\|_{L_{I}^{\infty}L_{x}^{\infty}} \leq \lambda^{\frac{a^{2}}{2} + am}\bar{B}_{0}, \label{eq:cG-132}
\end{equation}

\begin{equation}
\|\partial_{x}^{\frac{1}{2} + m}(\partial^{2}g_{\leq\lambda^{a}})\|_{L_{I}^{\infty}L_{x}^{2}} \leq \lambda^{\frac{a^{2}}{2} + am}\bar{B}_{0}, \label{eq:cG-133}
\end{equation}

\begin{equation}
\left\|\partial^{m}\square_{g_{\leq\lambda^{a}}}g_{\leq\lambda^{a}}\right\|_{L_{I}^{1}L_{x}^{\infty}} \leq \lambda^{-(1 - a) + am}\bar{B}_{0}.\label{eq:cG-134}
\end{equation}

常数 $\bar{B}_0$ 只依赖于常数 $M_0$ 与 $B_0$. 
\end{proposition}

\begin{remark*}
注意, 由构造, 度量 $g_{\leq\lambda^{a}}$ 的频率只在相对于空间变量 $x$ 的对偶 Fourier 变量的意义上被截断到 $\lambda^{a}$ 以上. 因此, 关于 $x$ 的每次微分至多引入一个 $\lambda^a$ 的因子. 然而, 利用 $g_{\leq \lambda^a}$ 依赖于波动方程的解这一事实, 我们可以对时间导数作出同样的结论. 
\end{remark*}

定理 (A3) 可以改写\footnote{因此我们可以完全忘记度量 $g_{\leq\lambda^{a}}$ 的来源, 我们只需要知道 (130)-(134). }为关于如下结果:在满足 \propref{prop:cG-6.14} 的估计的背景下度量 $g_{\leq\lambda^{a}}$ 的线性波动方程的解, 在某个固定小子区间 $I$ 上的局部 Strichartz 估计. 

\begin{theorem}[A4]\label{thm:cG-6.15}
设 $\psi$ 是线性波动方程

\begin{equation}
\begin{aligned}
&\square_{g_{\leq\lambda^{a}}}\psi = -\partial_{t}^{2}\psi + g_{\leq\lambda^{a}}^{ij}\partial_{i}\partial_{j}\psi = 0,  \\
&\psi|_{t = 0} = \psi_{0},  \quad \partial_{t}\psi|_{t = 0} = \psi_{1}
\end{aligned}\label{eq:cG-135}
\end{equation}

在长度 $|I|\leq \lambda^{-(1 - a)}$ 的时间区间 $I$ 上的解, 其初值数据 $\psi [0]$ 在 Fourier 空间中支集在集合 $\{\xi :\frac{1}{2}\lambda \leq |\xi |\leq 2\lambda \}$ 内. 假设度量 $g_{\lambda^a}$ 以参数 $a$ 满足 \propref{prop:cG-6.14} 的 \eqref{eq:cG-130}--\eqref{eq:cG-134}, 其中参数 $a$ 满足 $a < -1 + \sqrt{3}$. 设 $P_{\lambda}$ 是 Fourier 空间中集合 $\{\xi :\frac{1}{2}\lambda \leq |\xi |\leq 2\lambda \}$ 上的投影. 则对足够大的参数 $\Lambda$、所有二进 $\lambda \geq \Lambda$ 与固定的 $\epsilon >0$, 

\begin{equation}
\left\|P_{\lambda}\partial\psi\right\|_{L_{I}^{2}L_{x}^{\infty}} \leq C(\bar{B}_{0})|I|^{\epsilon}\|\psi[0]\|_{\dot{H}^{2 + \epsilon}}\label{eq:cG-136}
\end{equation}

其中常数 $C(\bar{B}_0)$ 与 $\lambda$ 无关. 
\end{theorem}

\mtsubsec{第 5 步：重标度}

把问题 \eqref{eq:cG-135} 换成它的重标度版本是方便的, 使得初值数据的支集具有频率 $|\xi| \sim 1$, 时间区间 $I$ 具有长度 $\leq \lambda^a$. 

必要时平移时间, 我们可以假设时间区间 $I$ 从 $t = 0$ 开始. 引入重标度度量族 $h_{\lambda}$

\begin{equation}
h_{\lambda}(t,  x) = g_{\leq\lambda^{a}}(\lambda^{-1}t,  \lambda^{-1}x)\label{eq:cG-137}
\end{equation}

\propref{prop:cG-6.14} 蕴含 $h_{\lambda}$ 在时间区间\footnote{我们保留记号 $I$ 表示重标度后的时间区间}$I = [0,  t_{*}]$($t_{*} \leq \lambda^{a}$)上满足如下估计\footnote{按照我们的约定, $A \lesssim B$ 表示对某个普适常数 $C$ 有 $A \leq C\cdot B$. 由自举假设, 我们估计中的所有常数都可能依赖于常数 $B_0$. 因此我们可以把 $B_0$ 看作普适常数, 并在下文中用 $\lesssim$ 替换对它的依赖. 此外, 如定理 (A3) 那样, 大频率 $\Lambda$ 的选择将由条件 $A \lesssim B$ 可替换为 $A \leq \Lambda^{\epsilon}B$(对任意正的 $\epsilon$)来决定. }:

\begin{equation}
\|\partial^{1 + m}h_{\lambda}\|_{L_{I_{\lambda}}^{1}L_{x}^{\infty}} \lesssim \lambda^{-(1 - a)(m + 1)}, \label{eq:cG-138}
\end{equation}

\begin{equation}
\left\|\partial^{1 + m}h_{\lambda}\right\|_{L_{I_{\lambda}}^{2}L_{x}^{\infty}} \lesssim \lambda^{-\frac{2 - a}{2} - (1 - a)m}, \label{eq:cG-139}
\end{equation}

\begin{equation}
\|\partial^{1 + m}h_{\lambda}\|_{L_{I_{\lambda}}^{\infty}L_{x}^{\infty}} \lesssim \lambda^{-1 + \frac{a^{2}}{2} - (1 - a)m}, \label{eq:cG-140}
\end{equation}

\begin{equation}
\left\|\partial_{x}^{\frac{1}{2} + m}\left(\partial^{2}h_{\lambda}\right)\right\|_{L_{I_{\lambda}}^{\infty}L_{x}^{2}} \lesssim \lambda^{-1 + \frac{a^{2}}{2} - (1 - a)m}, \label{eq:cG-141}
\end{equation}

\begin{equation}
\left\|\partial^{m}\bar{\Box}_{h_{\lambda}}h_{\lambda}\right\|_{L_{I_{\lambda}}^{1}L_{x}^{\infty}} \lesssim \lambda^{-(2 - a) - (1 - a)m}, \label{eq:cG-142}
\end{equation}

重标度之后, 定理 (A4) 变换为

\begin{theorem}[A5]\label{thm:cG-6.17}
设 $\psi$ 是线性波动方程

\begin{equation}
\begin{aligned}
&\Box_{h_{\lambda}}\psi = -\partial_{t}^{2}\psi + h_{\lambda}^{ij}\partial_{i}\partial_{j}\psi = 0,  \\
&\psi|_{t = 0} = \psi_{0},  \quad \partial_{t}\psi|_{t = 0} = \psi_{1}
\end{aligned}\label{eq:cG-143}
\end{equation}

在时间区间 $[0,  t_{*}]$($t_{*} \leq \lambda^{a}$)上的解. 假设参数 $\lambda \geq \Lambda$(对某个足够大的常数 $\Lambda$), 度量 $h_{\lambda}$ 以参数 $a$($a < -1 + \sqrt{3}$)满足 \eqref{eq:cG-138}--\eqref{eq:cG-142}. 设 $P$ 是 Fourier 空间中集合 $\{\xi : 1 \leq |\xi| \leq 2\}$ 上的投影算子. 则

\begin{equation}
\left\|P\partial\psi\right\|_{L_{[0,  t_{*}]}^{2}L_{x}^{\infty}} \lesssim |t_{*}|^{\epsilon}\big(\left\|\partial\psi_{0}\right\|_{L_{x}^{2}} + \left\|\psi_{1}\right\|_{L_{x}^{2}}\big)\label{eq:cG-144}
\end{equation}

其中不等式 $\lesssim$ 中的常数与 $\lambda$ 无关. 
\end{theorem}

\begin{remark*}
注意, 定理 (A5) 不包含关于初值数据 $\psi[0]$ 的 Fourier 支集的任何假设. 
\end{remark*}

\mtsubsec{第 6 步：衰减估计}

标准 $TT^{*}$ 型论证的一个变体, 见 {[}29{]}, 使我们能把 Strichartz 估计 \eqref{eq:cG-144} 化归为相应的色散不等式, 见 \eqref{eq:cG-145}. 在这个过程中, 我们用几何波动方程

\[
\Box_{h_{\lambda}}\psi = -\frac{1}{\sqrt{\det h_{\lambda}}}\partial_t\sqrt{\det h_\lambda}\partial_t\psi + \frac{1}{\sqrt{\det h_\lambda}}\partial_i(h_\lambda^{ij}\sqrt{\det h_\lambda}\partial_j\psi) = 0
\]

替换方程 $\Box_{h_{\lambda}}\psi = 0$\footnote{就 Strichartz 估计而言, 这两个波算子只差低阶项. }. 

\begin{theorem}[A6]\label{thm:cG-6.19}
设 $\psi$ 是线性波动方程

\begin{equation}
\begin{aligned}
&\Box_{h_{\lambda}}\psi = -\frac{1}{\sqrt{\det h_{\lambda}}}\partial_{t}\sqrt{\det h_{\lambda}}\partial_{t}\psi + \frac{1}{\sqrt{\det h_{\lambda}}}\partial_{i}(h_{\lambda}^{ij}\sqrt{\det h_{\lambda}}\partial_{j}\psi) = 0,  \\
&\psi|_{t = 0} = \psi_{0},  \quad \partial_{t}\psi|_{t = 0} = \psi_{1}
\end{aligned}\label{eq:cG-145}
\end{equation}

在时间区间 $[0,  t_{*}]$($t_{*} \leq \lambda^{a}$)上的解, 其初值数据 $\psi[0]$ 在 Fourier 空间中支集在集合 $\{\xi : \frac{1}{2} \leq |\xi| \leq 2\}$ 内. 我们只考虑参数 $\lambda \geq \Lambda$ 的大值. 假设度量 $h_{\lambda}$ 以参数 a($a < -1 + \sqrt{3}$)满足 \eqref{eq:cG-138}--\eqref{eq:cG-142}. 则对所有 $t \leq t_{*}$ 与某个固定、任意小的 $\epsilon > 0$

\begin{equation}
\|P\partial\psi(t)\|_{L_{x}^{\infty}} \lesssim \frac{1}{(1 + |t|)^{1 - \epsilon}}\|\psi[0]\|_{L_{x}^{1}}.\label{eq:cG-146}
\end{equation}
\end{theorem}

我们作最后的化简:把物理空间中的初值数据 $\psi[0]$ 分解为一系列本质上不交的支集之和, 这些支集包含在半径为 $\frac{1}{2}$ 的球内. 利用 $L^{1}$ 范数的可加性与标准 Sobolev 不等式, 我们可以把色散不等式 \eqref{eq:cG-146} 化归为一个 $L^{2}-L^{\infty}$ 衰减估计. 

\begin{theorem}[B]\label{thm:cG-6.20}
设 $\psi$ 是线性波动方程 \eqref{eq:cG-145} 在时间区间 $[0,  t_{*}]$($t_{*} \leq \lambda^{a}$)上的解, 其初值数据 $\psi[0]$ 在物理空间中支集在以原点为中心、半径为 $\frac{1}{2}$ 的球 $B_{\frac{1}{2}}(0)$ 内. 我们固定一个大的常数 $\Lambda$, 只考虑参数 $\lambda \geq \Lambda$ 的大值. 假设度量 $h_{\lambda}$ 以参数 $a$($a < -1 + \sqrt{3}$)满足 \eqref{eq:cG-138}--\eqref{eq:cG-142}. 则对所有 $t \leq t_{*}$、任意小的 $\epsilon > 0$ 与足够大的整数 $m > 0$, 

\begin{equation}
\|P\partial\psi(t)\|_{L_{x}^{\infty}} \lesssim \frac{1}{(1 + |t|)^{1 - \epsilon}}\sum_{k = 1}^{m}\|\partial^{k}\psi[0]\|_{L_{x}^{2}}.\label{eq:cG-147}
\end{equation}
\end{theorem}

\section*{参考文献}
\markboth{参考文献}{参考文献}
\addcontentsline{toc}{section}{参考文献}

\begin{mybibliography}{76}

\bibitem[1]{gfref1} H. Bahouri 与 J.Y. Chemin,  拟线性波动方程与色散效应,  \emph{Amer. J. Math.} 121 (1999),  1337--1377

\bibitem[2]{gfref2} H. Bahouri 与 J.Y. Chemin,  拟线性波动方程与 Strichartz 估计,  \emph{Int. Math. Res. Not.},  21 (1999),  1141--1177

\bibitem[3]{gfref3} M. Beals,  半线性波动方程解的自动扩展与奇点强度,  \emph{Ann. Math.} 118 (1983),  187--214

\bibitem[4]{gfref4} J. Bourgain,  某些格点子集的 傅里叶变换限制现象及其在非线性波动方程中的应用 I:Schrödinger 方程, II:KdV 方程,  \emph{Geom. Funct. Anal.} 3 (1993),  107--156,  209--262

\bibitem[5]{gfref5} N. Bournaveas,  三维空间中 Maxwell-Dirac 方程的局部存在性,  \emph{Comm. P.D.E.} 21 (1996),  5--6,  693--720

\bibitem[6]{gfref6} N. Bournaveas,  Dirac-Klein-Gordon 方程能量类解的局部存在性,  \emph{Comm. P.D.E.} 24 (1999),  7--8,  1167--1993

\bibitem[7]{gfref7} Y. Ch. Bruhat,  某些非线性偏微分方程组解的存在性定理,  \emph{Acta Math.} 88 (1952),  141--225

\bibitem[8]{gfref8} Y. Ch. Bruhat 与 R.P. Geroch,  广义相对论中 Cauchy 问题的整体方面,  \emph{Comm. Math. Phys.} 14,  329--335

\bibitem[9]{gfref9} L. Carlesson 与 P. Sjolin,  振荡积分与圆盘的乘子问题,  \emph{St. Math.} 44,  1972,  pp. 287--299

\bibitem[10]{gfref10} D. Christodoulou,  关于整体初值问题与奇点问题,  \emph{Class. Quant. Gr.} (1999) A23--A35

\bibitem[11]{gfref11} D. Christodoulou 与 A. Shadi Tahvildar-Zadeh,  关于球对称波映射的正则性,  \emph{Comm. Pure Appl. Math.} 46 (1993),  1041--1091

\bibitem[12]{gfref12} D. Christodoulou 与 S. Klainerman,  Minkowski 空间的整体非线性稳定性,  Princeton University Press (1993)

\bibitem[13]{gfref13} S. Cuccagna,  关于 $R^{3+1}$ 中 Maxwell-Klein-Gordon 系统的局部存在性,  \emph{Comm. PDE} 24 (1999),  no. 5--6,  851--867

\bibitem[14]{gfref14} D. Foschi 与 S. Klainerman,  波动方程的齐次 $L^{2}$ 双线性估计,  \emph{Ann. Scient. ENS} 4$^{e}$ série,  23 (2000),  211--274

\bibitem[15]{gfref15} F. Hélein,  从曲面到带对称性流形的弱调和映射的正则性,  \emph{Manusc. Math.} 70 (1991),  203--218

\bibitem[16]{gfref16} 非线性双曲微分方程讲义,  \emph{Mathématiques et Applications},  Springer-Verlag,  1996

\bibitem[17]{gfref17} F. John 与 S. Klainerman,  三维空间中非线性波动方程的近整体存在性,  \emph{Comm. Pure Appl. Math.} 36,  1980,  pp. 325--344

\bibitem[18]{gfref18} L. Kapitansky,  非线性波动方程的整体唯一弱解,  \emph{Math. Res. Lett.} 1 (1994),  211--223

\bibitem[19]{gfref19} M. Keel,  $R^{3+1}$ 中临界幂次 Yang-Mills-Higgs 的整体存在性,  \emph{Comm. PDE} 22 (1997),  1167--1227

\bibitem[20]{gfref20} M. Keel 与 T. Tao,  端点 Strichartz 估计,  \emph{Amer. J. Math.} 120 (1998),  no. 5,  955--980

\bibitem[21]{gfref21} M. Keel 与 T. Tao,  粗糙数据下 $R^{1+1}$ 上波映射的局部与整体适定性,  \emph{Int. Math. Res. Not.} 1998,  no. 21,  1117--1156

\bibitem[22]{gfref22} M. Keel 与 T. Tao,  能量范数以下 Maxwell-Klein-Gordon 方程大数据的整体适定性,  preprint

\bibitem[23]{gfref23} C. Kenig,  G. Ponce 与 L. Vega,  通过压缩原理得到的广义 KdV 方程的适定性与散射结果,  \emph{Comm. Pure Appl. Math.} 46 (1993),  no. 4,  527--620

\bibitem[24]{gfref24} C. Kenig,  G. Ponce 与 L. Vega,  负指标 Sobolev 空间中 KdV 方程的 Cauchy 问题,  \emph{Duke Math. J.} 71 (1994),  1--21

\bibitem[25]{gfref25} L. Kapitanski,  二阶线性双曲方程解在 Besov 与 Lizorkin-Triebel 空间中的范数估计,  \emph{J. Soviet Math.},  vol. 49 (1990),  pp. 1166--1186

\bibitem[26]{gfref26} S. Klainerman,  非线性波动方程解的长时间行为,  \emph{Proc. ICM} 1983,  Warszawa

\bibitem[27]{gfref27} S. Klainerman,  一致衰减估计与经典波动方程的 Lorentz 不变性,  \emph{Comm. Pure Appl. Math.} 35 (1985),  321--332

\bibitem[28]{gfref28} S. Klainerman,  零条件与非线性波动方程的整体存在性,  \emph{Lect. Appl. Math.} 23 (1986),  293--326

\bibitem[29]{gfref29} S. Klainerman,  PDE 作为统一的研究对象,  将发表于国际会议《数学中的愿景》(Tel Aviv 1999)论文集

\bibitem[30]{gfref30} S. Klainerman,  Strichartz 型不等式的交换向量场方法及其在拟线性波动方程中的应用,  将发表于 IMRN. S. Klainerman,  Strichartz 型不等式的交换向量场方法及其在拟线性波动方程中的应用,  preprint 2000,  将发表于 IMRN

\bibitem[31]{gfref31} S. Klainerman 与 M. Machedon,  零形式的时空估计与局部存在性定理,  \emph{Comm. Pure Appl. Math.},  46 (1993),  1221--1268

\bibitem[32]{gfref32} S. Klainerman 与 M. Machedon,  关于波动方程的正则性性质,  \emph{Physics on Manifolds},  M. Flato、R. Kerner 与 A. Lichnerowitz 编,  Kluwer Academic Publishers,  46 (1993),  1221--1268

\bibitem[33]{gfref33} S. Klainerman 与 M. Machedon,  关于有限能量的 Maxwell-Klein-Gordon 方程,  \emph{Duke Math. J.} 74 (1994),  19--44

\bibitem[34]{gfref34} S. Klainerman 与 M. Machedon,  $R^{3+1}$ 中 Yang-Mills 方程的有限能量解,  \emph{Ann. Math.} 142 (1995),  39--119

\bibitem[35]{gfref35} S. Klainerman 与 M. Machedon,  零形式的平滑化估计及其应用,  \emph{Duke Math. J.} 81 (1995),  99--133

\bibitem[36]{gfref36} S. Klainerman 与 M. Machedon,  关于与波映射相关的模型问题的正则性性质,  \emph{Duke Math. J.} 87 (1997),  no. 3,  553--589

\bibitem[37]{gfref37} S. Klainerman 与 M. Machedon,  关于 Strichartz 型不等式的注记,  \emph{Int. Math. Res. Not.},  no. 5 (1996),  201--220

\bibitem[38]{gfref38} S. Klainerman 与 M. Machedon,  零形式的估计与空间 $H_{s, \delta}$,  \emph{Int. Math. Res. Not.},  no. 17 (1996),  853--866

\bibitem[39]{gfref39} S. Klainerman 与 M. Machedon,  关于规范场理论的最优局部正则性,  \emph{Differential and Integral Equations} 10 (1997),  1019--1030

\bibitem[40]{gfref40} S. Klainerman 与 F. Nicolo,  关于广义相对论中 Cauchy 问题的局部与整体方面,  \emph{Class. Quant. Grav.} 16 (1999),  R73--R157. 还将出版一本同名的新书. 

\bibitem[41]{gfref41} S. Klainerman 与 I. Rodnianski,  关于临界 Sobolev 范数中波映射的整体正则性,  将发表于 IMRN (2001)

\bibitem[42]{gfref42} S. Klainerman 与 I. Rodnianski,  三维拟线性波动方程改进的局部适定性,  preprint 2001,  已投稿 Duke Math. Journ.

\bibitem[43]{gfref43} S. Klainerman,  I. Rodnianski 与 T. Tao,  Minkowski 空间中双线性估计的向量场方法,  撰写中

\bibitem[44]{gfref44} S. Klainerman 与 S. Selberg,  关于波映射型方程最优正则性的注记,  \emph{Comm. PDE} 22 (1997),  901--918

\bibitem[45]{gfref45} S. Klainerman 与 S. Selberg,  双线性估计及其在非线性波动方程中的应用,  已投稿 Notices of AMS

\bibitem[46]{gfref46} S. Klainerman 与 T. Sideris,  关于三维非相对论波动方程的近整体存在性,  \emph{Comm. Pure Appl. Math.} 49,  1996,  pp. 307--321

\bibitem[47]{gfref47} S. Klainerman 与 D. Tataru,  关于 $R^{4+1}$ 中 Yang-Mills 方程的最优局部正则性,  \emph{J. Amer. Math. Soc.} 12 (1999),  93--116

\bibitem[48]{gfref48} H. Lindblad,  关于小初值数据非线性波动方程解的寿命,  \emph{Comm. Pure Appl. Math.} 43,  1990,  pp. 445--472

\bibitem[49]{gfref49} H. Lindblad,  半线性波动方程局部存在性的反例,  \emph{Amer. J. Math.} 118 (1996),  1--16

\bibitem[50]{gfref50} H. Lindblad 与 C. Sogge,  关于半线性波动方程最小正则性的存在性与散射,  \emph{J. Funct. Anal.} 130 (1995),  357--526

\bibitem[51]{gfref51} M. Machedon,  零形式与非线性波动方程的 Fourier 分析,  \emph{Documenta Mathematica},  extra volume ICM,  1998,  III 49--55

\bibitem[52]{gfref52} G. Mockenhaupt,  A. Seeger 与 C. Sogge,  Fourier 积分算子的局部平滑化与 Carleson-Sjölin 估计,  \emph{J. Amer. Math. Soc.},  vol. 6 (1993),  pp. 65--130

\bibitem[53]{gfref53} Y. Meyer,  关于 J. M. Bony 的一个定理的注记,  \emph{Proceedings of the Seminar on Harmonic Analysis} (Pisa,  1980),  \emph{Rend. Circ. Mat. Palermo} (2) 1981,  suppl. 1,  1--20

\bibitem[54]{gfref54} G. Ponce 与 T. Sideris,  三维空间中非线性波动方程的局部正则性,  \emph{Comm. PDE} 18 (1993),  169--177

\bibitem[55]{gfref55} S. Selberg,  多重线性时空估计及其在非线性波动方程局部存在性理论中的应用,  Ph.D. 学位论文,  Princeton University 1999

\bibitem[56]{gfref56} S. Selberg,  关于非线性波动方程适定性的若干注记,  preprint 2000

\bibitem[57]{gfref57} S. Selberg,  波映射与双线性时空估计,  preprint 2000

\bibitem[58]{gfref58} J. Shatah 与 M. Struwe,  几何波动方程,  \emph{Courant Lecture Notes in Mathematics} 2 (1998)

\bibitem[59]{gfref59} J. Shatah 与 M. Struwe,  临界增长半线性波动方程在能量空间中的适定性,  \emph{Int. Math. Res. Not.} 7 (1994),  303--309

\bibitem[60]{gfref60} H. Smith,  具有 $C^{1, 1}$ 系数波动方程的参数层构造,  \emph{Annales de L'Institut Fourier},  vol. 48 (1998),  pp. 797--835

\bibitem[61]{gfref61} H. Smith 与 C. Sogge,  关于低正则性度量的 Strichartz 与特征函数估计,  \emph{Math. Res. Lett.},  vol. 1 (1994),  pp. 729--737

\bibitem[62]{gfref62} H. Smith 与 D. Tataru,  低正则性度量 Strichartz 估计的锐利反例,  Preprint

\bibitem[63]{gfref63} T. Sideris,  Minkowski 空间中调和映射的整体存在性,  \emph{Comm. Pure Appl. Math.} 42,  1989,  1--13

\bibitem[64]{gfref64} T. Sideris,  零条件与非线性波的整体存在性,  \emph{Invent. Math.} 123,  1996,  pp. 323--342

\bibitem[65]{gfref65} T. Sideris,  预应力非线性弹性波的非共振与整体存在性,  \emph{Annals of Math.} 151,  pp. 849--874

\bibitem[66]{gfref66} M. Struwe,  几何与数学物理中的演化问题,  \emph{AMS Prospects in Mathematics},  H. Rossi 编,  83--101

\bibitem[67]{gfref67} R.S. Strichartz,  傅里叶变换到二次曲面的限制与波动方程解的衰减,  \emph{Duke Math. J.} 44 (1977),  705--714

\bibitem[68]{gfref68} T. Tao,  锥的端点双线性限制定理与若干锐利零形式估计,  preprint,  已投稿 Math. Z.

\bibitem[69]{gfref69} T. Tao,  $L^2$ 函数的多重线性加权卷积及其在非线性色散方程中的应用,  preprint

\bibitem[70]{gfref70} T. Tao,  波映射的整体正则性 I,  将发表于 I.M.R.N.

\bibitem[71]{gfref71} T. Tao,  波映射的整体正则性 II,  已投稿 C.M.P.

\bibitem[72]{gfref72} D. Tataru,  波映射的局部与整体结果 I,  \emph{Comm. PDE} 23 (1998),  1781--1793

\bibitem[73]{gfref73} D. Tataru,  关于波映射方程的整体存在性与散射 II,  将发表于 AJM

\bibitem[74]{gfref74} D. Tataru,  非光滑系数算子的 Strichartz 估计与非线性波动方程,  将发表于 Amer. J. Math.

\bibitem[75]{gfref75} D. Tataru,  非光滑系数二阶双曲算子的 Strichartz 估计 III,  preprint

\bibitem[76]{gfref76} T. Wolff,  锐利双线性锥限制估计,  将发表于 Ann. Math.

\end{mybibliography}
